% ============================================================
% Chapter 2: Background and Related Work
%
% PURPOSE:
%   Establish what other researchers and practitioners have done
%   to address your problem and related problems. Show what your
%   work builds on and what gap it fills.
%
% THIS CHAPTER GOES BY DIFFERENT NAMES depending on your thesis:
%   - "Literature Review" — for research/survey-focused theses
%   - "Background and Related Work" — for systems/implementation
%   - "Background" — if you need a separate technical primer
%   Rename the \chapter{} title below to match yours.
%
% SUGGESTED STRUCTURE:
%   Option A (research thesis):
%     2.1 — [First thematic area]
%     2.2 — [Second thematic area]
%     2.3 — [Third thematic area]
%     2.4 — Summary and Identified Gap
%
%   Option B (systems/implementation thesis):
%     2.1 — Relevant Technical Background
%     2.2 — Existing Tools and Approaches
%     2.3 — Comparison and Gap
%
% TIPS:
%   - Synthesize sources — do not summarize papers one by one.
%   - Group related work by theme, not by author.
%   - End with a section that explicitly identifies the gap
%     your thesis addresses.
%   - Cite frequently: \cite{KEY} produces [1] in the output.
%
% TARGET LENGTH: 8--15 pages.
% TARGET CITATIONS: 15--25 sources minimum.
% ============================================================

\chapter{Background and Related Work}
% Rename to "Literature Review" or "Background" if appropriate.
\label{chapter2}

\noindent
[Write a brief introductory paragraph describing how this chapter
is organized and what it covers. Example: "This chapter reviews
work relevant to [your topic], organized into three themes:
[Theme A], [Theme B], and [Theme C]. The chapter concludes by
identifying the gap this thesis addresses."]


\section{[First Theme or Technical Area]}

\noindent
[Describe the first major area of related work. Group multiple
sources together and synthesize their contributions rather than
summarizing each paper individually. What do these works agree on?
Where do they differ? What are their limitations
\cite{example_journal, example_conference}?]


\section{[Second Theme or Technical Area]}

\noindent
[Describe the second major area \cite{example_book}. Add or remove section blocks to match the structure of your literature.]


\section{[Third Theme or Technical Area]}
% Delete this section if you only have two themes.
% Add more sections with \section{} if you have more.

\noindent
[Describe the third major area \cite{example_arxiv}.]


\section{Summary and Identified Gap}
% This section is important regardless of your thesis type.
% Explicitly name what existing work has not addressed.

\noindent
[Summarize what the reviewed work has established and where it
falls short. Connect the gap directly to your research questions
from Chapter 1. Be specific — "no existing work addresses X in
the context of Y" is more convincing than "more work is needed."]
